\begin{helper}
Fix \(I\subseteq[n]\), integers \(k,\ell_R,\ell_B\), and parameters
\(m,p,\varepsilon\).  On the projection-size event
\[
S=\{\omega:
\noderef{TwoBiteProjectionSizeEvent}(I,k,\ell_R,\ell_B,\omega)\}
\]
there is a finite disjoint partition \(S=\dot\bigcup_{h\in H}C_h\) with the
following adaptive cylinder structure.  For every history \(h\) there are a
representative sample \(\omega_h\in C_h\), a recorded coordinate set \(R_h\),
a terminal candidate set \(U_h\), and a terminal order \(L_h\) listing
\(U_h\) without repetition, such that
\[
\omega\in C_h
\quad\Longleftrightarrow\quad
\omega\hbox{ has the same injection as }\omega_h
\hbox{ and agrees with }\omega_h
\hbox{ on every coordinate of }R_h.
\tag{1}
\]
The recorded set \(R_h\) is the pre-terminal recorded set:
\[
e\in R_h
\quad\Longleftrightarrow\quad
e\in\noderef{TwoBitePreTerminalRecordedPairs}(\omega,\varepsilon,I)
\qquad(\omega\in C_h).
\tag{2}
\]
In particular, every coordinate of \(R_h\) is in the canonical terminal
coordinate universe: it is tagged red or blue, has distinct endpoints, and is
written in increasing endpoint order.
The terminal set is the full unrecorded terminal coordinate universe:
\[
e\in U_h
\quad\Longleftrightarrow\quad
e\in\noderef{TwoBiteTerminalCoordinateUniverse}(m)
\hbox{ and }e\notin R_h. \tag{3}
\]
Thus membership in \(C_h\) fixes only the injection and the \(R_h\)-truth
table.  It imposes no condition on any coordinate of \(U_h\).

For every \(\omega\in C_h\),
\[
e\in\noderef{TwoBiteStagedOpenPairs}(\omega,\varepsilon,I)
\quad\Longrightarrow\quad
e\in U_h. \tag{4}
\]
The reverse implication is not asserted: \(U_h\) may contain coordinates that
are not projected pairs of \(I\), are touched, or become \(C^+(I)\)
same-color closed during the terminal scan.

The order \(L_h\) is prefix safe.  If
\[
L_h=P,e,Q,
\]
and \(\omega\in C_h\), then any second sample \(\omega'\) with the same
injection as \(\omega\), the same edge values as \(\omega\) on \(R_h\), and
the same edge values as \(\omega\) on the prefix \(P\), has the same terminal
decision data as \(\omega\) for the current coordinate \(e\):
\[
e\in\noderef{TwoBiteStagedOpenPairs}(-,\varepsilon,I),\qquad
\noderef{TwoBiteProjectionPairTouched}(-,\varepsilon,I,e),\qquad
\noderef{TwoBiteProjectionPairSameColorClosed}(-,I,e),\qquad
e\in\noderef{TwoBitePreTerminalRecordedPairs}(-,\varepsilon,I).
\tag{5}
\]
No agreement is required on the current coordinate \(e\) or on any later
coordinate of \(Q\).
\end{helper}
\begin{proof}
Apply \(\noderef{FixedSetPreTerminalHistoryPartition}\) to
\(I,k,\ell_R,\ell_B,m,p,\varepsilon\).  Denote the finite history type by
\(H\), its cells by \(C_h\), its recorded sets by \(R_h\), and its
representatives by \(\omega_h\).  We record the child facts with the names
used below; the recorded-coordinate hygiene (8'') is then derived from (8')
and the definition of the pre-terminal recorded set:
\[
\omega\in S \Longleftrightarrow \exists h\in H,\ \omega\in C_h, \tag{5}
\]
\[
h\ne h'\Longrightarrow C_h\cap C_{h'}=\varnothing, \tag{6}
\]
\[
\omega_h\in C_h, \tag{7}
\]
\[
\omega\in C_h
\Longleftrightarrow
\omega\hbox{ has the same injection as }\omega_h
\hbox{ and agrees with }\omega_h\hbox{ on }R_h, \tag{8}
\]
\[
\omega\in C_h
\Longrightarrow
\bigl(e\in R_h\Longleftrightarrow
e\in\noderef{TwoBitePreTerminalRecordedPairs}(\omega,\varepsilon,I)\bigr),
\tag{8'}
\]
\[
e\in R_h\Longrightarrow e\hbox{ is written with increasing endpoints},
\tag{8''}
\]
\[
\omega\in C_h\Longrightarrow \omega\in S, \tag{9}
\]
\[
\omega\in C_h,\ \omega'\equiv\omega\hbox{ on the injection and on }R_h
\Longrightarrow
\begin{cases}
\noderef{TwoBiteStageRevealedVertex}(\omega,\varepsilon,I,x)
\Longleftrightarrow
\noderef{TwoBiteStageRevealedVertex}(\omega',\varepsilon,I,x),\\
\noderef{TwoBiteProjectionPairTouched}(\omega,\varepsilon,I,c)
\Longleftrightarrow
\noderef{TwoBiteProjectionPairTouched}(\omega',\varepsilon,I,c),
\end{cases}
\tag{10}
\]
for every base vertex \(x\) and every tagged coordinate \(c\), and finally
\[
\omega\in C_h,\ e\in\noderef{TwoBiteStagedOpenPairs}
(\omega,\varepsilon,I)
\Longrightarrow
e\in\noderef{TwoBiteTerminalCoordinateUniverse}(m)
\hbox{ and }e\notin R_h. \tag{11}
\]
The disjointness in (6) is the canonical pre-terminal transcript
disjointness proved in that child: the injection and the \(R_h\)-truth table
uniquely replay the same pre-terminal recorded set and table.

For each fixed history \(h\), apply
\(\noderef{FixedSetTerminalSupportClassification}\) to the recorded set
\(R_h\).  This supplies a terminal set \(U_h\) and an order \(L_h\).  Its
conclusions are
\[
e\in U_h
\Longleftrightarrow
e\in\noderef{TwoBiteTerminalCoordinateUniverse}(m)\hbox{ and }e\notin R_h,
\tag{12}
\]
\[
L_h\hbox{ has no repetitions and }(L_h)_{\mathrm{finset}}=U_h, \tag{13}
\]
\[
e\in U_h\Longrightarrow e\notin R_h, \tag{14}
\]
\[
e\in U_h\Longrightarrow e\hbox{ is written with increasing endpoints},
\tag{15}
\]
and the support-classification property: if
\[
L_h=P,e,Q,
\tag{16}
\]
then every one-sided coordinate that can support a
\(C^+(I)\) same-color witness for \(e\) lies in
\[
R_h\cup P_{\mathrm{finset}}. \tag{17}
\]
Concretely, for \(e=\operatorname{red}(a,b)\) and every red base vertex
\(r<a,b\), both \(\operatorname{red}(r,a)\) and
\(\operatorname{red}(r,b)\) lie in (17), and the blue case is the identical
blue-copy assertion.  The choices of \(U_h\) and \(L_h\) for all histories
are legitimate because the history type \(H\) is finite.

We now verify the conjuncts of the statement in order.  The finite partition,
the cover of \(S\), the disjointness of distinct cells, the representative
condition, the exact cylinder description, the exact pre-terminal recorded-set
identity, and the containment \(C_h\subseteq S\) are precisely (5), (6), (7),
(8), (8'), and (9).  The order
conjunct
\[
L_h.\mathrm{Nodup}\quad\hbox{and}\quad L_h.\mathrm{toFinset}=U_h
\]
is (13).  The recorded-coordinate hygiene (8'') is obtained as follows:
given \(e\in R_h\), apply (8') with the representative
\(\omega_h\in C_h\).  Then
\[
e\in\noderef{TwoBitePreTerminalRecordedPairs}
(\omega_h,\varepsilon,I).
\]
By the definition of \(\noderef{TwoBitePreTerminalRecordedPairs}\), this
finite set is a filter of
\(\noderef{TwoBiteTerminalCoordinateUniverse}(m)\).  Hence \(e\) lies in the
terminal coordinate universe, and the definition of that universe says
exactly that \(e\) is a tagged red or blue non-loop coordinate in increasing
endpoint order.  The terminal-universe identity (3) is (12), the statement
that terminal coordinates are outside \(R_h\) is (14), and the increasing
non-loop coordinate hygiene is (15).  Since (8) fixes only the injection and
the \(R_h\)-truth table, and since (12) makes \(U_h\) disjoint from \(R_h\),
no cell condition fixes any terminal coordinate value.  Moreover, because
both recorded and terminal coordinates are canonical, a terminal coordinate
cannot be fixed indirectly by a reversed recorded coordinate.

Next prove the staged-open inclusion (4).  Let \(\omega\in C_h\) and
suppose
\[
e\in\noderef{TwoBiteStagedOpenPairs}(\omega,\varepsilon,I).
\]
By (11), \(e\) lies in
\(\noderef{TwoBiteTerminalCoordinateUniverse}(m)\) and does not lie in
\(R_h\).  Substituting these two facts into (12) gives \(e\in U_h\).  This is
only the asserted one-way inclusion; \(U_h\) may still contain coordinates
that are not projected pairs of \(I\), are touched, or are
\(C^+(I)\) same-color closed.

It remains to prove prefix safety.  Fix \(h\), \(\omega\in C_h\), and a
decomposition \(L_h=P,e,Q\).  Let \(\omega'\) have the same injection as
\(\omega\), agree with \(\omega\) on every coordinate of \(R_h\), and agree
with \(\omega\) on every coordinate in \(P_{\mathrm{finset}}\).  Since
\(\omega\in C_h\), the cylinder identity (8) says that \(\omega\) has the
same injection as \(\omega_h\) and agrees with \(\omega_h\) on \(R_h\).  The
same-injection assumption between \(\omega\) and \(\omega'\), combined with
the same-injection relation between \(\omega\) and \(\omega_h\), gives the
same injection relation between \(\omega'\) and \(\omega_h\).  Likewise, if
\(c\in R_h\), then \(\omega'\) and \(\omega\) have the same truth value for
\(\noderef{TwoBiteEdgeCoordinateValue}(-,c)\), and \(\omega\) and
\(\omega_h\) have the same truth value for that coordinate.  Chaining these
two equivalences gives agreement between \(\omega'\) and \(\omega_h\) on
every coordinate of \(R_h\).  Applying (8) in the reverse direction therefore
puts \(\omega'\) in the same history cell \(C_h\).  This is the only point at
which membership of \(\omega'\) in \(C_h\) is used, and it is obtained without
any information about the current coordinate \(e\) or the suffix \(Q\).

Applying the pre-terminal replay clause (10) to \(\omega\) and \(\omega'\) gives, for
every base vertex \(x\),
\[
\noderef{TwoBiteStageRevealedVertex}(\omega,\varepsilon,I,x)
\Longleftrightarrow
\noderef{TwoBiteStageRevealedVertex}(\omega',\varepsilon,I,x),
\tag{18}
\]
and, for every tagged coordinate \(c\),
\[
\noderef{TwoBiteProjectionPairTouched}(\omega,\varepsilon,I,c)
\Longleftrightarrow
\noderef{TwoBiteProjectionPairTouched}(\omega',\varepsilon,I,c).
\tag{19}
\]
The common injection also fixes the red and blue projection images of \(I\).
Indeed, unfolding the red and blue projections shows that the projected image
of each final vertex of \(I\) is computed from the common injection alone;
hence the two red image sets and the two blue image sets are equal.  Unfolding
\(\noderef{TwoBiteProjectionPairSet}\), which is the tagged union of base
pairs in those two image sets, gives for every coordinate \(c\)
\[
c\in\noderef{TwoBiteProjectionPairSet}(\omega,I)
\Longleftrightarrow
c\in\noderef{TwoBiteProjectionPairSet}(\omega',I).
\tag{19a}
\]
In particular, projected-pair membership for the current coordinate \(e\) is
fixed before any edge value at \(e\) is queried.

The exact-recorded identity (8') applies to both \(\omega\) and \(\omega'\),
because both samples lie in \(C_h\).  Hence, for the current coordinate \(e\),
\[
e\in\noderef{TwoBitePreTerminalRecordedPairs}(\omega,\varepsilon,I)
\Longleftrightarrow e\in R_h
\Longleftrightarrow
e\in\noderef{TwoBitePreTerminalRecordedPairs}(\omega',\varepsilon,I).
\tag{19'}
\]

For same-color closedness of \(e\), take the support set
\[
S_h(P)=R_h\cup P_{\mathrm{finset}}.
\]
The displayed decomposition of \(L_h\) lets us use (17).  In the red case
\(e=\operatorname{red}(a,b)\), every possible same-color witness for
\(\noderef{TwoBiteProjectionPairSameColorClosed}(-,I,e)\) uses, for some red
base vertex \(r<a,b\), only the two one-sided coordinates
\(\operatorname{red}(r,a)\) and \(\operatorname{red}(r,b)\); (17) places both
coordinates in \(S_h(P)\).  In the blue case, every possible blue witness is
supported on \(\operatorname{blue}(s,a)\) and \(\operatorname{blue}(s,b)\),
and (17) again places both coordinates in \(S_h(P)\).  The samples agree on
all of \(S_h(P)\): if a support coordinate lies in \(R_h\), use the recorded
agreement; if it lies in \(P_{\mathrm{finset}}\), use the prefix agreement.
These two cases cover the whole support set by its definition, and neither
case mentions the current coordinate \(e\) unless \(e\) was already recorded
or in the prefix.  For the displayed split \(L_h=P,e,Q\), that cannot happen:
\(L_h\) has no repetitions, \(e\in U_h\) because it appears in the list
listed by \(L_h\), and \(U_h\) is disjoint from \(R_h\).  The argument
therefore never asks for a fresh value of \(e\), and it never asks for any
coordinate in \(Q\).
Thus \(\noderef{FixedSetSameColorClosedWitnessCylinder}\), applied with
support \(S_h(P)\), gives
\[
\noderef{TwoBiteProjectionPairSameColorClosed}(\omega,I,e)
\Longleftrightarrow
\noderef{TwoBiteProjectionPairSameColorClosed}(\omega',I,e).
\tag{20}
\]

Finally unfold the definition of
\(\noderef{TwoBiteStagedOpenPairs}\) at the current coordinate \(e\).  The
projected-pair predicate is fixed by (19a), the touched
predicate is fixed by (19) with \(c=e\), and the same-color-closed predicate
is fixed by (20).  The pre-terminal recorded predicate is fixed by (19').
Hence
\[
e\in\noderef{TwoBiteStagedOpenPairs}(\omega,\varepsilon,I)
\Longleftrightarrow
e\in\noderef{TwoBiteStagedOpenPairs}(\omega',\varepsilon,I).
\]
Equation (19) with \(c=e\) gives the touched equivalence, and (20) gives the
same-color-closed equivalence; (19') gives the pre-terminal recorded
equivalence.  The argument used the common injection, recorded agreement on
\(R_h\), and prefix agreement on \(P_{\mathrm{finset}}\) only.  It did not use
the value of the current coordinate \(e\), and it did not use any coordinate
in the suffix \(Q\).  This proves the full adaptive cylinder package.
\end{proof}
